%%%%%%%%%%%%%%%%%%%%%%%%%%%%%%%%%
%
%       EXERCÍCIO
%
%%%%%%%%%%%%%%%%%%%%%%%%%%%%%%%%%

\ifdefstring{\atividade}{prova}{%
    \ifdefstring{\modo}{objetivo}{%
        \renewcommand{\valorquestao}{\ValorQObj\ ponto}
    }{%
        \renewcommand{\valorquestao}{\ValorQDisc\ pontos}
    }%
}%

\begin{exercicioBanco}[\valorquestao]
Considerando um intervalo de tempo de 10 anos a partir de hoje, o valor de uma máquina deprecia linearmente com o tempo, isto é, o valor da máquina \(y\), em função do tempo \(x\), é dado por uma função afim \(y = ax + b\). Se o valor da máquina daqui a dois anos for R\$ \(6400{,}00\), e seu valor daqui a cinco anos e meio for R\$ \(4300{,}00\), determine qual será o seu valor daqui a sete anos.

% Define as alternativas
\newcommand{\alternativas}{%
    \begin{center}
        \begin{tabularx}{\textwidth}{XXXXX}
            (a) R\$ \(2800{,}00\). &
            (b) R\$ \(3000{,}00\). &
            (c) R\$ \(3200{,}00\). &
            (d) R\$ \(3400{,}00\). &
            (e) R\$ \(3600{,}00\).
        \end{tabularx}
    \end{center}
}

% Resposta correta
\newcommand{\resposta}{D}

% Lógica condicional para exibição
\ifdefstring{\atividade}{lista}{%
    \alternativas
    \vspace{0.5em}
    
    \noindent\textbf{Resposta:} letra \textbf{\resposta}.
}{%
    \ifdefstring{\modo}{objetivo}{%
        \alternativas
    }{%
        % modo = discursiva → não mostra alternativas
    }
}
\end{exercicioBanco}

